\section{Expected Results}

The understanding of material behavior under realistic operating conditions is essential for advancing catalytic and electrochemical energy conversion technologies. In systems such as SOFC/SOEC cells, and catalytic reactors for steam reforming of ethanol (SRE), the working principles are governed by mechanisms involving surface catalysis, electrocatalytic charge transfer, defect chemistry, and ionic and electronic transport in complex oxides. These processes are highly sensitive to temperature, gas composition, and electrochemical processes. And they frequently induce dynamic structural and electronic changes that cannot be captured by ex-situ characterization. To address this challenge, the present project emphasizes the development of advanced instrumentation based on synchrotron techniques capable of probing materials in-situ and operando.

The experimental strategy will be implemented at CNPEM using the fourth-generation synchrotron light source Sirius. In particular, the QUATI beamline will be employed due to its capability for in-situ and operando material characterization. The high brilliance, stability, and energy tunability of the source enable high-resolution X-ray absorption spectroscopy, providing element-specific information on oxidation state, coordination environment, and local atomic structure under realistic reaction conditions. Custom-designed high-temperature electrochemical and catalytic cells will be integrated into the beamline environment, allowing measurements under controlled gas atmospheres and applied electrical bias. This approach ensures that catalytic, electrocatalytic, and transport mechanisms are measured while the material is in working conditions, thereby establishing direct correlations between atomic-scale dynamics and macroscopic performance.

These experiments are crucial for elucidating degradation mechanisms that limit device durability. Materials employed in SOFC/SOEC devices and ethanol reforming catalysts are subjected to redox cycling, high thermal loads, chemical gradients, and reactive intermediates that promote phase segregation, ion migration, carbon deposition, and microstructural degradation. Operando X-ray absorption spectroscopy, combined with electrochemical measurements, enables early detection of such phenomena by monitoring changes in electronic structure, local bonding, and electro-catalytic efficiency. 

The results will significantly advance the understanding of the processes, and support the optimization of compositions, microstructures, and operating conditions. Most importantly, the advanced instrumentation and standardized experimental setups developed within this project will be made available to the LNLS user community, expanding the analytical capabilities for a broad range of electrochemical devices and catalytic systems and fostering collaborative research across disciplines.

In summary, we expect to achieve significant results in the following areas; the development and optimization of new nano materials for SOFCs and SOECs. These materials are intended to be cost-effective and durable, thereby contributing to efficient energy generation and hydrogen production; the development and optimization of new catalysts for the steam reforming of ethanol, which is particularly important for enhancing the ethanol potential and deliver economic benefits; and the engineering of specialized, innovative cells for in-situ and operando investigations of solid oxide cells and catalytic processes.
